\documentclass[11pt]{article}

\usepackage[margin=25mm]{geometry}

\usepackage{amsmath,amssymb,fontspec,unicode-math}

\IfFontExistsTF{DejaVu Serif}{\setmainfont{DejaVu Serif}}{\setmainfont{Latin Modern Roman}}

\IfFontExistsTF{Latin Modern Math}{\setmathfont{Latin Modern Math}}{\setmathfont{DejaVu Math TeX Gyre}}

\usepackage{xurl,hyperref}

\setlength{\parindent}{0pt}\setlength{\parskip}{6pt}

\begin{document}

\title{Solver-Free Decision-Focused Learning for Linear Optimization Problems}

\author{ELVIA preliminary AI analysis}\date{21 October 2026}\maketitle

Preliminary assessment: Minor.\par

Independent scientific review remains pending. Author adoption has not been requested.

Three release-level discrepancies affect neighbor masking, relaxed-solution data flow and runtime measurement. The adjacent-vertex theorem survives, with a separate zero-margin tie qualification.

\section*{Review scope}

Full supplied technical report read alongside the corresponding original-paper claim passages and applicable author-release source. This pass independently checks the stated reasoning and bounds; historical executions were not rerun.

\section*{Problem 1: Value-based padding deletes a real zero neighbor}

Source: Eq.(4), PDF p.4; methods/lava.py, LAVA.forward, lines 13–21; dataset\_augmented.py, slack projection

Assessment: supported; impact: minor

\subsection*{Source claim (paraphrase)}

LAVA sums truncated adjacent-vertex objective gaps. The release treats an all-zero projected vertex as padding.

\subsection*{Analysis}

Consider minimization on \(0\leq x\leq1\), represented by \(x+s=1\), with true cost \(-x\). The true optimum is \((1,0)\), whose genuine neighbor \((0,1)\) projects to \(x=0\) after slack removal. Predicted cost \(\hat c=1\) selects \(x=0\) with true regret \(1\). The real neighbor gap is \(1\), but the value-based padding test deletes it, yielding zero loss and zero gradient. Separately, mapping gaps below \(-\varepsilon\) to zero differs from the printed clamp: for \(\varepsilon=1/10\), gaps \(-1/5,-1/20\) violate scalar convexity since the midpoint maps to \(0>-1/40\). Their ordinary derivatives nevertheless agree away from the threshold.

\subsection*{Finding and scope}

The zero-neighbor mask changes the defined surrogate and its gradients on an allowed projected LP. The threshold-to-zero scalar difference alone does not establish different updates away from the kink. Proposition 4.1 remains intact.

\subsection*{Proposed correction}

Carry an explicit neighbor-validity mask or count independent of coordinates, and clamp real gaps at \(-\varepsilon\) to match Eq.(4).

The described correction remains a proposal; author adoption or a corresponding erratum is not established.

\section*{Problem 2: Integer-solution cache and filtering corrupt relaxed neighborhoods}

Source: Section 5.3, PDF p.8; data/dataset\_augmented.py, lines 72–90

Assessment: supported; impact: minor

\subsection*{Source claim (paraphrase)}

The knapsack workflow trains against relaxed optima, but caches neighbors by the integer optimum and removes neighbors equal to that integer solution.

\subsection*{Analysis}

For \(3x+2y\leq4\), \(x,y\in\{0,1\}\), use costs \((4,3)\) then \((5,3)\). Both integer optima are \(I=(1,0)\); relaxed optima are respectively \(A=(2/3,1)\) and \(B=(1,1/2)\). Their neighbor sets are \(N(A)=\{B,(0,1)\}\) and \(N(B)=\{I,A\}\). Reusing the cache under key \(I\) supplies stale neighbors at \(B\). With predicted maximization cost \(q=(-1,-3)\), stale gaps are \(0,-1/2\), so the release returns no update, while the omitted actual neighbor \(I\) has gap \(3/2\). Even without a cache collision, filtering by the integer solution deletes \(I\). The selected integer decision is the origin, with true regret \(5\) under the second cost.

\subsection*{Finding and scope}

Both data-flow defects are concrete for the supported knapsack/relaxation interface. Their occurrence and effect in the original benchmark arrays are not established.

\subsection*{Proposed correction}

Key by the actual relaxed vertex and constraint identity, or disable caching. Remove a self-neighbor only against that same relaxed label; retain all genuine adjacent vertices.

The described correction remains a proposal; author adoption or a corresponding erratum is not established.

\section*{Problem 3: Training-time accumulator excludes most training work}

Source: Section 5.1, PDF pp.7–8; utils/train\_model.py, lines 43–44 and 66–103

Assessment: supported; impact: minor

\subsection*{Source claim (paraphrase)}

The paper describes training time including adjacent-vertex precomputation and actual training, excluding validation, and uses a \(600\)-second budget. The release times only the loss evaluation.

\subsection*{Analysis}

The prediction the regressor forward call precedes the timer; elapsed time is added before zero\_grad, backward and optimizer.step. Hence an iteration records approximately \(t_{\rm loss}\), rather than \(t_{\rm pred}+t_{\rm loss}+t_{\rm back}+t_{\rm opt}\). These omitted operations are training, not validation. Precomputation is separately included.

\subsection*{Finding and scope}

The release timing definition does not match the stated protocol. It cannot recover a corrected speedup or establish that the archived commit generated the published plots.

\subsection*{Proposed correction}

Time the complete prediction, loss, backward and optimizer step with a monotonic clock; add precomputation and consistently exclude validation. Remeasure the affected runtime comparisons if this path generated them.

The described correction remains a proposal; author adoption or a corresponding erratum is not established.

\section*{Problem 4: Zero-margin optimality leaves solver ties unresolved}

Source: Proposition 4.1 and sentence preceding Eq.(4), PDF p.4

Assessment: not\_issue; impact: minor

\subsection*{Source claim (paraphrase)}

Making a labeled vertex predicted-cost optimal is described as giving zero regret, whereas the regret evaluates one selected optimizer.

\subsection*{Analysis}

For the unit square with true cost \((1,1)\), true optimum \((0,0)\), and nonzero predicted cost \((1,0)\), both \((0,0)\) and \((0,1)\) are predicted optima. The adjacent-vertex zero-margin loss is zero, but a solver may choose \((0,1)\), with true regret \(1\). Positive strict adjacent margins instead give a unique optimum, and the experiments use \(\varepsilon=0.1\).

\subsection*{Finding and scope}

This is a tie-handling qualification of a zero-margin interpretation, not a refutation of Proposition 4.1 or the positive-margin method. Tie-independent zero regret requires every predicted optimizer to be a true optimizer.

\section*{Limits}

Preliminary AI source/argument assessment; independent human scientific adjudication remains pending.\par

No author-provided or retained verification program was executed. Published empirical results were not independently reproduced.\par

\section*{Source identity}

NeurIPS 2025 official proceedings PDF (exact retained local copy inspected); compared with current authors’ arXiv:2505.22224v2

\url{https://proceedings.neurips.cc/paper_files/paper/2025/file/e7dd11514cc84c83622e4f5766f9eafa-Paper-Conference.pdf}

Reviewed PDF SHA-256: \nolinkurl{a9b1f09d2eb9e4145529779586758c48eafb27eed7cf9403f52c3681f4bddafb}

The SHA-256 identifies the exact retained PDF inspected in this rapid source/argument check. Original supplied report bytes are retained locally. The original PDF retrieval timestamp is taken from matching retained metadata.

\end{document}
